\documentclass[11pt,a4paper]{article}
\usepackage[utf8]{inputenc}
\usepackage[T1]{fontenc}
\usepackage[spanish]{babel}
\usepackage[margin=2.5cm]{geometry}
\usepackage{csquotes}
\usepackage[backend=biber,style=ieee]{biblatex}
\usepackage{hyperref}
\usepackage{enumitem}
\addbibresource{refs.bib}

\newcommand{\campo}[1]{\par\medskip\noindent\textbf{#1}\par\smallskip}
\newcommand{\pendiente}{\textbf{[PENDIENTE]}}

\begin{document}
	
	% ---------------- Caratula ----------------
	\begin{center}
		{\Large\textbf{Universidad T\'ecnica Estatal de Quevedo}}\\[4pt]
		Facultad de Ciencias de la Computaci\'on\\
		Carrera de Software\\[12pt]
		\textbf{Aplicaciones Distribuidas (ISR-701)}\\
		\textbf{FORM-FB-01: Fichas bibliogr\'aficas}\\
		Fundamentos de sistemas distribuidos y arquitecturas distribuidas\\[12pt]
		Estudiante: Reinoso V\'elez Eduardo David\\
		Nivel y paralelo: S\'eptimo nivel, paralelo A\\
		Fecha: 6 de noviembre de 2026\\
		Repositorio: \url{https://github.com/ereinosov/app-dist-form-fb-01}
	\end{center}
	
	\bigskip
	\noindent\textbf{Sistema elegido para TA-IND-01:} Apache Cassandra, comparado con Amazon DynamoDB.
	
	\section*{Ficha 1. Ferreira et al. (2024): niveles de consistencia en NoSQL}
	\campo{Referencia (IEEE)}
	\fullcite{ferreira2024impacts}
	\campo{DOI verificado}
	10.1186/s13677-024-00716-7
	\campo{Resumen (100 palabras)}
	Ferreira y su equipo quieren saber cuánto rendimiento se pierde al exigir más consistencia a una base NoSQL (Not only SQL, no relacional). Prueban Cassandra, MongoDB y Redis en Google Cloud, cada una por separado, con la herramienta de carga Locust y entre 1 y 10 usuarios simultáneos. Cassandra corre con tres nodos y factor de replicación 3, y comparan el nivel ONE (basta la respuesta de una réplica) con ALL (responden todas). Lo repiten con los nodos en una región y en varias. El resultado es claro: la consistencia fuerte baja el rendimiento y alarga las respuestas, y la distancia entre regiones agrava el efecto.
	\campo{Dato concreto (con p\'agina o secci\'on)}
	En la Tabla 8 (p. 14), pasar de consistencia débil a fuerte en Cassandra con réplicas en varias regiones reduce el rendimiento de escritura un 97,54 \% y sube el tiempo de respuesta un 6.028,34 \%. En una sola región, las lecturas pierden un 62,30 \% de rendimiento y tardan un 175,15 \% más.
	\campo{Nota personal}
	Esta fuente me da números concretos para la parte de CAP (consistencia, disponibilidad y tolerancia a particiones) y PACELC (si hay partición, disponibilidad o consistencia; si no, latencia o consistencia) de TA-IND-01. Tengo presente que su nivel fuerte es ALL y que solo un nodo recibió las peticiones, así que no puedo usar sus cifras como si fueran de QUORUM (mayoría de réplicas).
	
	\newpage \section*{Ficha 2. Rehman, Aguiar y Barraca (2022): tolerancia a fallos en la nube}
	\campo{Referencia (IEEE)}
	\fullcite{rehman2022fault}
	\campo{DOI verificado}
	10.1109/ACCESS.2022.3182211
	\campo{Resumen (100 palabras)}
	Rehman, Aguiar y Barraca revisan más de cien trabajos sobre tolerancia a fallos en la nube. Empiezan fijando tres términos: falta (la causa), error (el estado que se desvía de lo esperado) y fallo (cuando el sistema deja de cumplir su función). Después reparten las técnicas en reactivas, que actúan cuando el problema ya ocurrió, y proactivas, que intentan adelantarse. Relacionan métricas de componente y de sistema, comparan arquitecturas existentes y cierran con temas abiertos: nube, niebla y borde, aprendizaje automático y consumo de energía. Su conclusión es que combinar ambos enfoques probablemente rinda más que elegir uno solo.
	\campo{Dato concreto (con p\'agina o secci\'on)}
	Secciones IV.F y IV.G. Entre las técnicas reactivas citan checkpoint/restart (guardar el estado y reiniciar desde ahí), replicación, migración de tareas y reintento. Entre las proactivas, autorreparación, migración preventiva, rejuvenecimiento de software y balanceo de carga. La sección VII advierte que las proactivas pueden costar más por el monitoreo continuo.
	\campo{Nota personal}
	Me sirve de marco para ordenar lo que hacen Cassandra y DynamoDB ante un fallo: la replicación y los reintentos quedan del lado reactivo, y con esta clasificación puedo decirlo con las mismas palabras en la sección de tolerancia a fallos de TA-IND-01.
	
	\section*{Ficha 3. Kumari y Kaur (2021): revisi\'on de tolerancia a fallos}
	\campo{Referencia (IEEE)}
	\fullcite{kumari2021survey}
	\campo{DOI verificado}
	10.1016/j.jksuci.2018.09.021
	\campo{Resumen (100 palabras)}
	Kumari y Kaur escriben una revisión sobre cómo mantener un servicio en la nube cuando algo falla. Primero describen la infraestructura y los modelos de despliegue y de servicio, y luego repasan los enfoques que ya existían en sistemas distribuidos. El núcleo es la clasificación en técnicas proactivas y reactivas, con ejemplos como la replicación, el reintento y la migración, y un apartado sobre su uso en aplicaciones basadas en la nube. Proponen una tabla de métricas para evaluarlas, con fiabilidad, disponibilidad, tiempo de respuesta, rendimiento y costo. Cierran con líneas de investigación futura, entre ellas el aprendizaje profundo.
	\campo{Dato concreto (con p\'agina o secci\'on)}
	Secciones 4.2 y 4.3 y Tabla 3. Proactivas: autorreparación, rejuvenecimiento de software, migración preventiva y balanceo de carga. Reactivas: checkpointing, migración de tareas, replicación, reintento y reenvío de tareas. La Tabla 3 reúne las métricas para evaluar la tolerancia a fallos: rendimiento, tiempo de respuesta, fiabilidad, disponibilidad y costo, entre otras.
	\campo{Nota personal}
	La clasificación coincide en gran parte con la de la ficha 2, pero su lista de métricas (Tabla 3) es lo más aprovechable: la usaré para decidir con qué criterios comparar Cassandra y DynamoDB en TA-IND-01.
	
	\section*{Ficha 4. Huang et al. (2022): fallos metaestables}
	\campo{Referencia (IEEE)}
	\fullcite{huang2022metastable}
	\campo{ISBN verificado}
	978-1-939133-28-1 (actas de OSDI 22; el art\'iculo no tiene DOI)
	\campo{Resumen (100 palabras)}
	Huang y sus coautores estudian los fallos metaestables: una sobrecarga que sigue aunque ya haya desaparecido lo que la provocó. Para ver cómo ocurren en la práctica, revisan informes públicos de incidentes de empresas grandes y pequeñas. Separan el disparador (un pico de carga o una pérdida de capacidad) de lo que mantiene el problema, como los reintentos o la recolección de basura (liberación automática de memoria). Reproducen tres casos en aplicaciones de ejemplo de código abierto. Concluyen que estos fallos aparecen en las caídas más graves y que los reintentos son el factor de mantenimiento más frecuente.
	\campo{Dato concreto (con p\'agina o secci\'on)}
	Sección 2.2 (p. 75): los reintentos son el efecto que más veces mantiene el fallo, presente en más de la mitad de los incidentes. Resumen (p. 73): al menos 4 de las 15 caídas mayores de AWS (Amazon Web Services) en la última década fueron metaestables.
	\campo{Nota personal}
	Me da un caso real de fallo para la sección de tolerancia a fallos de TA-IND-01 y una advertencia útil: los reintentos de los clientes pueden empeorar la sobrecarga de un clúster, algo que tendré en cuenta al hablar de Cassandra bajo carga.
	
	\section*{Ficha 5. Elhemali et al. (2022): Amazon DynamoDB}
	\campo{Referencia (IEEE)}
	\fullcite{elhemali2022dynamodb}
	\campo{ISBN verificado}
	978-1-939133-29-8 (actas de USENIX ATC 22; el art\'iculo no tiene DOI)
	\campo{Resumen (100 palabras)}
	Elhemali y su equipo de AWS cuentan cómo operan DynamoDB a gran escala y cómo cambió su arquitectura desde 2012. Presentan seis propiedades del servicio: gestión completa, multiinquilino, escala sin límite, rendimiento predecible, alta disponibilidad y uso flexible. Cada partición se copia en varias zonas de disponibilidad y elige líder con Multi-Paxos (protocolo de consenso). Solo el líder atiende escrituras y lecturas fuertes. El texto explica también cómo manejan fallos silenciosos, particiones calientes y despliegues, y da cifras reales, como un pico de 89,2 millones de peticiones por segundo en el Prime Day de 2021.
	\campo{Dato concreto (con p\'agina o secci\'on)}
	p. 1040: una escritura se confirma cuando un quórum (mayoría) de réplicas guarda el registro; solo el líder atiende escrituras y lecturas fuertemente consistentes, mientras que cualquier réplica puede atender lecturas eventualmente consistentes. p. 1045: un quórum de escritura sano son 2 de 3 réplicas en zonas de disponibilidad distintas.
	\campo{Nota personal}
	Es la base de lo que voy a contar de DynamoDB en la comparación de TA-IND-01. Su esquema de un líder por partición contrasta con Cassandra, que no tiene nodo líder, y ahí se explica buena parte de la diferencia en consistencia.
	
	\printbibliography[title={Referencias}]
	
	\section*{Declaraci\'on de uso de inteligencia artificial}
	\textbf{Herramienta:} Claude (Anthropic), modelo Claude Sonnet 5.5.
	
	\textbf{Para qu\'e la us\'e:}
	(1) buscar en la web fuentes candidatas sobre consistencia, tolerancia a fallos y comunicaci\'on
	en sistemas distribuidos, con Apache Cassandra como sistema elegido para TA-IND-01;
	(2) generar la plantilla \LaTeX{} de este documento y el esqueleto del archivo \texttt{refs.bib}, y
	(3) cotejar los datos bibliogr\'aficos de las candidatas con p\'aginas de editoriales y
	repositorios.
	
	\textbf{C\'omo verifiqu\'e lo obtenido:} abr\'i personalmente el DOI o el enlace de cada una de las
	cinco fuentes y compar\'e autores, a\~no, t\'itulo, revista o libro, volumen, n\'umero y p\'aginas con
	\texttt{refs.bib}; correg\'i las diferencias; consult\'e en Scopus/SJR la indexaci\'on y el
	cuartil de cada revista; y le\'i cada fuente para localizar el dato concreto y su p\'agina.
	
	\textbf{Redacci\'on de las fichas:} la herramienta me ayud\'o a redactar los borradores iniciales de los res\'umenes y
	las notas, a partir de los res\'umenes de las fuentes. Yo verifiqu\'e cada afirmaci\'on, ubiqu\'e los datos concretos con su p\'agina y reescrib\'i los textos con
	mis propias palabras. Los datos y sus p\'aginas son los que constat\'e en las fuentes.
	
\end{document}